\Needspace{10\baselineskip}
\section{数据预处理}\label{sec_preprocessing}
为使四问模型使用一致的输入口径，本节统一整理附件1的烘房温湿环境和附件2的药材半径数据，将离散观测转换为数值求解所需的时间函数，并明确长期计算中的环境延续范围。

\subsection{数据整理与单位统一}
附件1给出时间、烘房温度和水分浓度，共241组观测，覆盖0至14400 s，间隔60 s。附件2给出时间与药材半径，共145组观测，覆盖0至259200 s，即72 h，间隔1800 s，即0.5 h。两份附件分别用于构造表面环境输入$T_a(t),C_a(t)$和第四问的几何输入$R(t)$。

计算中时间统一使用s，半径由cm换算为m；正文表格按题目要求再换算为h和cm。温度观测保留摄氏度，扩散系数经验式中的绝对温度由药材局部温度加273.15得到。附件1中的烘房温度只作为边界输入，不代替药材内部温度。水分量保留题给kg/kg口径，其气固交换解释见模型假设。

问题一截取附件1前1800 s的31组观测，问题二截取前10800 s的181组观测，问题三、四使用前14400 s的全部241组观测。各问均保留相应区间的原始观测和小幅波动，不作全局平滑；环境插值区间与本问的计算口径保持一致。

\subsection{烘房环境插值与时间延续}
数值求解器需要在观测时刻之间计算表面环境，因此对各问所需观测区间采用保形分段三次Hermite插值（PCHIP）~\cite{pchip}。记相邻观测点为$(t_j,f_j)$，区间长度为$h_j=t_{j+1}-t_j$，$z=(t-t_j)/h_j$，则
\begin{equation}\label{eq_pchip}
\begin{gathered}
f(t)=(2z^3-3z^2+1)f_j+(-2z^3+3z^2)f_{j+1}\\
+(z^3-2z^2+z)h_jm_j+(z^3-z^2)h_jm_{j+1},
\end{gathered}
\end{equation}
其中$m_j$由保形规则确定，$f$分别取烘房温度$T_a$和水分量$C_a$。式\eqref{eq_pchip}精确通过观测点，所得时间函数在求解器需要的时刻提供边界输入。前30 min的环境观测插值见图\ref{fig_input}，烘房升温不意味着药材内部同步升温。
\figone{q1_oven_input.pdf}{.76}{前30 min烘房温度和水分量随时间变化。曲线由附件1观测插值得到，输入升温不等于药材内部同步升温。}{fig_input}

附件1没有覆盖完整干燥时段。其最后1 h的温度和水分量均值分别为约$49.999\ ^\circ\mathrm C$和$0.04999$ kg/kg，已接近平台。根据假设（6），问题三、四在$t\le14400$ s时保留全部观测输入，随后采用基于末段平台值的环境延续假设
\begin{equation}\label{eq_late}
T_a(t)=50\ ^\circ\mathrm C,\qquad C_a(t)=0.05\ \mathrm{kg/kg},\qquad t>14400\ \mathrm s.
\end{equation}
式\eqref{eq_late}用于补全长期计算的边界条件，不是后期实测数据，也不是由PCHIP直接外推得到。切换时刻保留原观测终值，随后允许边界输入发生小幅变化，药材内部温湿状态保持连续。该延续假设的影响通过问题三的情景敏感性分析量化，并作为长期时长预测的适用限制保留。

\subsection{半径轨迹插值与数据覆盖范围}
附件2的半径由2 cm降至1.198 cm。将其换算为m后，采用与式\eqref{eq_pchip}相同的插值方法构造
\begin{equation}\label{eq_radius_input}
R(t)=\mathrm{PCHIP}\{(t_j,R_j)\},\qquad 0\le t\le259200\ \mathrm s.
\end{equation}
图\ref{fig_q4R}给出半径观测散点与插值曲线，插值精确通过观测点。$R(t)$仅在附件2的覆盖区间内使用，72 h是数据覆盖终点，不是预设的烘干答案；第四问计算在该范围内达标，无需对半径作数据外推。
\figone{q4_R_history.pdf}{.78}{附件2半径观测及PCHIP插值。半径由2 cm降至1.198 cm；72 h为观测覆盖终点，不是预设烘干时长。}{fig_q4R}

半径数据只确定外表面轨迹，不能唯一识别内部材料点的位移。第四问将以此输入建立比例径向收缩假设，并推导材料运动和热湿传递方程；这些模型关系不由数据插值本身给出。
\FloatBarrier
